\section{Evaluation Metrics}

The autoencoder models were evaluated using complementary metrics targeting
different aspects of reconstruction quality. Root mean squared error (RMSE)
was used to quantify pointwise agreement between the observed and
reconstructed BOLD time series. Functional connectivity (FC) similarity was
used as a secondary measure of preservation of static inter-regional
relationships, while sliding-window functional connectivity dynamics
(SwFCD) similarity was used as the primary evaluation metric. This
prioritisation reflects the main objective of the thesis: preserving the
dynamic organisation of BOLD activity rather than optimising pointwise
signal reconstruction alone.

\subsection{Time-Series Reconstruction Error}

Pointwise reconstruction accuracy was quantified using root mean squared
error (RMSE). For an observed BOLD signal
$\mathbf{X}\in\mathbb{R}^{T\times R}$ and its reconstruction
$\hat{\mathbf{X}}$, RMSE is defined as

\begin{equation}
    \operatorname{RMSE}
    =
    \sqrt{
        \frac{1}{TR}
        \sum_{t=1}^{T}
        \sum_{r=1}^{R}
        \left(
            X_{t,r}-\hat{X}_{t,r}
        \right)^2
    }.
\end{equation}

Lower RMSE values indicate greater pointwise agreement between the observed
and reconstructed BOLD signals. RMSE was used primarily to assess how
accurately the amplitude of the regional time series was reproduced.
However, pointwise reconstruction error alone does not determine whether
the functional relationships between regions, or their temporal variation,
are preserved. RMSE was therefore interpreted alongside the connectivity-
based evaluation metrics.

\subsection{Static Functional Connectivity Similarity}

Static functional connectivity was calculated from the Pearson correlation
between each pair of regional BOLD time series, producing a symmetric
functional connectivity matrix
$\mathbf{C}\in\mathbb{R}^{R\times R}$. Pearson correlation is widely used
to quantify statistical dependence between regional BOLD signals in
resting-state fMRI and forms a common basis for functional connectivity
analysis~\cite{Hutchison2013}.

To compare the observed and reconstructed FC matrices, only the upper
triangular elements were extracted, excluding the diagonal. This avoids
including duplicate values resulting from matrix symmetry and removes the
trivial self-correlations along the diagonal. Let
$\operatorname{ut}(\mathbf{C})$ denote this vectorisation operation. FC
similarity was then calculated as

\begin{equation}
    r_{\mathrm{FC}}
    =
    \operatorname{corr}
    \left(
        \operatorname{ut}\left(\mathbf{C}(\mathbf{X})\right),
        \operatorname{ut}\left(\mathbf{C}(\hat{\mathbf{X}})\right)
    \right).
\end{equation}

Higher values of $r_{\mathrm{FC}}$ indicate greater agreement between the
static inter-regional connectivity patterns of the observed and
reconstructed signals. FC correlation was treated as a secondary
evaluation metric because it describes the average functional organisation
over the complete scan but does not explicitly characterise changes in
connectivity over time.

\subsection{Dynamic Functional Connectivity Similarity}

Dynamic functional connectivity was evaluated using sliding-window
functional connectivity dynamics (SwFCD). Sliding-window approaches estimate
functional connectivity repeatedly over partially overlapping temporal
segments and are commonly used to characterise temporal variability in
resting-state functional connectivity~\cite{Hutchison2013}. SwFCD
extends this representation by comparing the similarity of window-specific
FC patterns across time, thereby describing the recurrence and temporal
organisation of functional connectivity states~\cite{Hansen2015}.

The BOLD signal was divided into overlapping windows of $W=30$ time points,
with a step size of $S=3$ time points. For each window, an FC matrix was
calculated and vectorised. Let
$\mathbf{F}_1,\ldots,\mathbf{F}_K$ denote the resulting window-specific FC
vectors. The SwFCD matrix was defined by the Pearson correlation between
each pair of window-specific FC patterns,

\begin{equation}
    \operatorname{SwFCD}_{i,j}
    =
    \operatorname{corr}
    \left(
        \mathbf{F}_i,
        \mathbf{F}_j
    \right),
\end{equation}

where $i$ and $j$ index the sliding windows.

As with the static FC metric, only the upper triangular portion of the
SwFCD matrix was used during evaluation. In addition, a region surrounding
the main diagonal was excluded before vectorisation. Consecutive
sliding windows overlap substantially and therefore share a large number
of time points. As a result, pairs of nearby windows naturally exhibit high
similarity, producing a band of elevated SwFCD values around the main
diagonal. Retaining this region could artificially increase the correlation
between observed and reconstructed SwFCD matrices even when their
longer-range dynamic structure differs.

The excluded diagonal offset was defined as

\begin{equation}
    k
    =
    \frac{W}{S}.
\end{equation}

For the values used in this work,
$W=30$ and $S=3$, giving an offset of

\begin{equation}
    k = 10.
\end{equation}

Only upper-triangular SwFCD elements outside this near-diagonal region were
therefore included in the evaluation. Let
$\operatorname{ut}_{k}(\cdot)$ denote this masked upper-triangular
vectorisation. SwFCD similarity was calculated as

\begin{equation}
    r_{\mathrm{SwFCD}}
    =
    \operatorname{corr}
    \left(
        \operatorname{ut}_{k}
        \left(
            \operatorname{SwFCD}(\mathbf{X})
        \right),
        \operatorname{ut}_{k}
        \left(
            \operatorname{SwFCD}(\hat{\mathbf{X}})
        \right)
    \right).
\end{equation}

Higher values of $r_{\mathrm{SwFCD}}$ indicate greater agreement between the
temporal organisation of functional connectivity in the observed and
reconstructed BOLD signals. SwFCD correlation was used as the primary
model-selection metric because preserving dynamic functional activity is
the central objective of this work. The use of a fixed window length, step
size, and masking procedure across all models ensured that SwFCD comparisons
were performed consistently. Sliding-window analyses are sensitive to
methodological choices such as window length and overlap, making such
consistency important for interpretation~\cite{Hutchison2013}.

\subsection{Model Comparison and Selection}

Candidate model configurations were evaluated across multiple random seeds.
Where more than one metric value was available for a particular seed, these
values were first averaged to obtain a single value for each model, seed,
and evaluation metric. Statistical comparisons were therefore performed at
the seed level.

Because corresponding model configurations were evaluated using matched
random seeds, differences between models were assessed using paired
$t$-tests. For two models $A$ and $B$, the paired differences for a metric
$m$ were defined as

\begin{equation}
    d_s
    =
    m_{A,s}
    -
    m_{B,s},
\end{equation}

where $s$ indexes the random seed. The paired $t$-statistic is

\begin{equation}
    t
    =
    \frac{\bar{d}}
    {s_d/\sqrt{n}},
\end{equation}

where $\bar{d}$ is the mean paired difference, $s_d$ is the standard
deviation of the paired differences, and $n$ is the number of seeds.

The significance threshold was set to
$\alpha=0.05$. Where multiple pairwise comparisons were performed, the
resulting $p$-values were adjusted using the Holm sequentially rejective
procedure~\cite{Holm1979}. The Holm procedure controls the family-wise
error rate while sequentially adjusting the significance criteria applied
to multiple hypotheses~\cite{Holm1979}.

Paired tests were performed for RMSE, FC correlation, and SwFCD correlation.
The statistical results were interpreted together with the corresponding
mean performance and seed-level distributions. Raincloud plots were used
throughout the experimental analysis to visualise the distribution of
results across seeds, together with their mean and variability. Means were
reported in the accompanying discussion to describe the direction and
magnitude of performance differences, while variability was discussed
explicitly where it materially affected interpretation.

Model selection considered all three evaluation metrics, with greatest
importance placed on SwFCD correlation because of the focus on preservation
of dynamic activity. FC correlation provided complementary information
about static connectivity preservation, while RMSE described pointwise
time-series reconstruction quality. Statistical significance alone was not
used to determine which model was preferable; the direction and magnitude
of the metric differences were considered jointly with the corresponding
Holm-adjusted $p$-values.

Where candidate configurations showed comparable overall performance and
no statistically significant improvement was identified, preference was
given to the computationally simpler model. Model simplicity was considered
in terms of factors such as architectural depth, number of trainable
parameters, and computational requirements. This criterion was used to
avoid introducing additional model complexity where it did not provide a
statistically supported improvement in reconstruction performance.
